\item \textbf{Cálculo de la temperatura exterior tocando la flauta.}
Es un flautista en una orquesta local. En un día frío de invierno, llega tarde a una actuación. Al llegar a la sala de orquesta, sabrá que ha perdido el ensayo del grupo antes de la presentación, por lo que afina su instrumento en el aire frío que se encuentra fuera de la puerta del escenario. Después de afinar, corre dentro del auditorio, donde la temperatura es de $22.2^\circ\text{C}$, toma asiento y comienza a tocar la primera canción con el resto de la orquesta.

Le da mucha vergüenza darse cuenta de que está tocando la canción medio paso más arriba que sus colegas de la orquesta. Su entusiasmo por la física supera su vergüenza musical al darse cuenta de que puede usar esta información para calcular la temperatura exterior. (Suponga que la longitud del instrumento no cambia con la temperatura. Un medio paso representa una relación de frecuencia de $2^{1/12}$).